\documentclass[11pt]{article}
\usepackage{amsmath,amssymb,amsthm}

\newtheorem{theorem}{Theorem}

\begin{document}

\title{Completeness of the Full Three-Equation System}
\author{}
\date{}
\maketitle

\paragraph{System.}
Consider
\[
a^2 + \left(\frac{a^2-r_2^2}{2r_2}\right)^2
=
\left(\frac{a^2+r_2^2}{2r_2}\right)^2,
\]
\[
a^2 + \left(\frac{a^2-r_3^2}{2r_3}\right)^2
=
\left(\frac{a^2+r_3^2}{2r_3}\right)^2,
\]
\[
\left(\frac{a^{2}-r_{2}^{2}}{2r_{2}}\right)^{2}
+
\left(\frac{a^{2}-r_{3}^{2}}{2r_{3}}\right)^{2}
= g^{2}.
\]
Define
\[
b_2:=\frac{a^2-r_2^2}{2r_2},\quad
c_2:=\frac{a^2+r_2^2}{2r_2},\quad
b_3:=\frac{a^2-r_3^2}{2r_3},\quad
c_3:=\frac{a^2+r_3^2}{2r_3}.
\]
Then the system is exactly
\[
a^2+b_2^2=c_2^2,\qquad a^2+b_3^2=c_3^2,\qquad b_2^2+b_3^2=g^2.
\]

\begin{theorem}[Rational completeness]
Let $a,b_2,c_2,b_3,c_3,g\in\mathbb Q$ satisfy
\[
a^2+b_2^2=c_2^2,\qquad
a^2+b_3^2=c_3^2,\qquad
b_2^2+b_3^2=g^2,
\]
with orientation assumptions
\[
b_2<c_2,\qquad b_3<c_3,
\]
and $g\neq 0$. Then there exist
\[
r_2,r_3,K\in\mathbb Q,\qquad m,n\in\mathbb Z,
\]
such that
\[
r_2\neq 0,\quad r_3\neq 0,
\]
\[
b_2=\frac{a^2-r_2^2}{2r_2},\quad
c_2=\frac{a^2+r_2^2}{2r_2},\quad
b_3=\frac{a^2-r_3^2}{2r_3},\quad
c_3=\frac{a^2+r_3^2}{2r_3},
\]
and simultaneously
\[
b_2=K(m^2-n^2),\qquad
b_3=2Kmn,\qquad
g=K(m^2+n^2).
\]
\end{theorem}

\begin{proof}
Split the argument in two standard completeness statements.

\medskip
\noindent\textbf{Step 1: first two equations.}
From
\[
a^2+b_i^2=c_i^2,\qquad b_i<c_i\qquad (i=2,3),
\]
set
\[
r_i:=c_i-b_i\neq 0.
\]
Then
\[
a^2=(c_i-b_i)(c_i+b_i)=r_i(c_i+b_i),
\]
so $c_i+b_i=a^2/r_i$. Solving
\[
c_i-b_i=r_i,\qquad c_i+b_i=\frac{a^2}{r_i}
\]
gives
\[
b_i=\frac{a^2-r_i^2}{2r_i},\qquad
c_i=\frac{a^2+r_i^2}{2r_i}.
\]
Hence the divisor parameterization is complete for the first two equations.

\medskip
\noindent\textbf{Step 2: third equation.}
Now use
\[
b_2^2+b_3^2=g^2,\qquad g\neq 0.
\]
By the rational Euclidean completeness theorem for Pythagorean triples, there exist
$K\in\mathbb Q$ and $m,n\in\mathbb Z$ such that
\[
b_2=K(m^2-n^2),\qquad
b_3=2Kmn,\qquad
g=K(m^2+n^2).
\]
Combining Step 1 and Step 2 gives the stated simultaneous parameterization.
\end{proof}

\begin{theorem}[Integer completeness (as a corollary)]
Let $a,b_2,c_2,b_3,c_3,g\in\mathbb Z$ satisfy the same three equations and
\[
b_2<c_2,\qquad b_3<c_3,\qquad g\neq 0.
\]
Then the same conclusion holds after embedding into $\mathbb Q$:
there exist $r_2,r_3,K\in\mathbb Q$ and $m,n\in\mathbb Z$ satisfying the same formulas.
\end{theorem}

\begin{proof}
Cast the integer equations and inequalities to $\mathbb Q$ and apply the rational theorem.
\end{proof}

\paragraph{What ``encompasses all solutions'' means here.}
Under the mild orientation convention $b_i<c_i$ (to avoid the $r_i=0$ denominator)
and the nondegenerate condition $g\neq 0$ for the Euclidean step, every rational or
integer solution of the full three-equation system is captured by the displayed
parameter families $(r_2,r_3)$ and $(K,m,n)$.

\end{document}
\documentclass[11pt]{article}
\usepackage{amsmath,amssymb,amsthm}

\newtheorem{theorem}{Theorem}
\newtheorem{proposition}[theorem]{Proposition}

\begin{document}

\title{Completeness of the Full Three-Equation System}
\author{}
\date{}
\maketitle

\paragraph{System.}
For $a,r_2,r_3,g\in\mathbb{Q}$ with $r_2r_3\neq 0$, consider
\[
a^2+\left(\frac{a^2-r_2^2}{2r_2}\right)^2
=
\left(\frac{a^2+r_2^2}{2r_2}\right)^2,
\]
\[
a^2+\left(\frac{a^2-r_3^2}{2r_3}\right)^2
=
\left(\frac{a^2+r_3^2}{2r_3}\right)^2,
\]
\[
\left(\frac{a^2-r_2^2}{2r_2}\right)^2
+
\left(\frac{a^2-r_3^2}{2r_3}\right)^2
=
g^2.
\]

Define
\[
b_2:=\frac{a^2-r_2^2}{2r_2},\quad
c_2:=\frac{a^2+r_2^2}{2r_2},\quad
b_3:=\frac{a^2-r_3^2}{2r_3},\quad
c_3:=\frac{a^2+r_3^2}{2r_3}.
\]
Then the system is exactly
\[
a^2+b_2^2=c_2^2,\qquad
a^2+b_3^2=c_3^2,\qquad
b_2^2+b_3^2=g^2.
\]

\begin{proposition}[Soundness]
For every $a,r_2,r_3,g\in\mathbb{Q}$ with $r_2r_3\neq 0$, if
\[
\left(\frac{a^2-r_2^2}{2r_2}\right)^2
+
\left(\frac{a^2-r_3^2}{2r_3}\right)^2
=
g^2,
\]
then $(a,b_2,c_2,b_3,c_3,g)$ satisfies
\[
a^2+b_2^2=c_2^2,\qquad
a^2+b_3^2=c_3^2,\qquad
b_2^2+b_3^2=g^2.
\]
\end{proposition}

\begin{proof}
The first two identities are polynomial identities (clear denominators and expand):
\[
4a^2r_i^2+(a^2-r_i^2)^2=(a^2+r_i^2)^2,\qquad i\in\{2,3\}.
\]
So $a^2+b_i^2=c_i^2$ for $i=2,3$. The third equality is exactly the displayed hypothesis after substituting the definitions of $b_2,b_3$.
\end{proof}

\begin{theorem}[Completeness over $\mathbb{Q}$]
Let $a,b_2,c_2,b_3,c_3,g\in\mathbb{Q}$ satisfy
\[
a^2+b_2^2=c_2^2,\qquad
a^2+b_3^2=c_3^2,\qquad
b_2^2+b_3^2=g^2,
\]
and assume
\[
b_2<c_2,\qquad b_3<c_3.
\]
Then there exist nonzero $r_2,r_3\in\mathbb{Q}$ such that
\[
b_2=\frac{a^2-r_2^2}{2r_2},\qquad
c_2=\frac{a^2+r_2^2}{2r_2},\qquad
b_3=\frac{a^2-r_3^2}{2r_3},\qquad
c_3=\frac{a^2+r_3^2}{2r_3},
\]
and
\[
\left(\frac{a^2-r_2^2}{2r_2}\right)^2
+
\left(\frac{a^2-r_3^2}{2r_3}\right)^2
=
g^2.
\]
One may take
\[
r_2=c_2-b_2,\qquad r_3=c_3-b_3.
\]
\end{theorem}

\begin{proof}
From the two Pythagorean equations and $b_i<c_i$, the standard divisor reconstruction gives
\[
r_i:=c_i-b_i\neq 0,\quad
b_i=\frac{a^2-r_i^2}{2r_i},\quad
c_i=\frac{a^2+r_i^2}{2r_i}\qquad (i=2,3).
\]
Now substitute these expressions into $b_2^2+b_3^2=g^2$; this yields
\[
\left(\frac{a^2-r_2^2}{2r_2}\right)^2
+
\left(\frac{a^2-r_3^2}{2r_3}\right)^2
=
g^2.
\]
Hence every rational solution of the three-equation system (with the orientation convention $b_i<c_i$) is encompassed by this parameterization.
\end{proof}

\begin{theorem}[Integer-input corollary]
If $a,b_2,c_2,b_3,c_3,g\in\mathbb{Z}$ satisfy
\[
a^2+b_2^2=c_2^2,\qquad
a^2+b_3^2=c_3^2,\qquad
b_2^2+b_3^2=g^2,
\]
with $b_2<c_2$ and $b_3<c_3$, then there exist nonzero $r_2,r_3\in\mathbb{Q}$ satisfying the same formulas above. Thus the same divisor system is complete for integer solutions (with rational parameters).
\end{theorem}

\begin{proof}
View the integer data in $\mathbb{Q}$ and apply the previous theorem.
\end{proof}

\paragraph{Conclusion.}
Under the same orientation assumptions used in the Lean development, the displayed divisor formulas (plus the third equation for $g$) capture all rational solutions and therefore all integer solutions.

\end{document}
